\documentclass[10pt]{article}
\usepackage[margin=0.78in]{geometry}
\usepackage{amsmath,amssymb,amsthm,mathrsfs}
\usepackage{microtype}
\usepackage{xcolor}
\usepackage{hyperref}
\hypersetup{colorlinks=true,linkcolor=blue!55!black,urlcolor=blue!55!black,citecolor=blue!55!black}
\newtheorem{theorem}{Theorem}
\newtheorem{lemma}[theorem]{Lemma}
\newcommand{\R}{\mathbb R}
\newcommand{\cS}{\mathcal S}
\newcommand{\Jac}{\mathscr J}
\title{A factor-two obstruction to the global-phase\newline Jacobian variational problem}
\author{Candidate full-solution packet for arXiv:1611.00653}
\date{11 August 2026}

\begin{document}
\maketitle

\begin{center}
\fbox{\parbox{0.92\linewidth}{\textbf{Status:} candidate full negative answer, likely valid.  The Schwartz-class variational construction requested in the Discussion of Section~7.2 cannot exist.  For every measurable set $E$, its proposed Jacobian integral is at most one half of the total Dirichlet energy.}}
\end{center}

\section{The question}

Carbonaro and Dragi\v cevi\'c ask whether there is a measurable
$E\subset\R^2$, with both $E$ and $E^c$ of positive measure, such that for
every $\delta>0$ one can find real functions $\rho>0$ and $\omega$ with
$u=\rho e^{i\omega}\in\cS(\R^2)$ and
\begin{equation}\label{eq:asked}
 \int_{\R^2}|\nabla u|^2
 <(1+\delta)\int_E\Jac(\rho^2,\omega).
\end{equation}
Here $\Jac(F,G)=F_xG_y-F_yG_x$.  This is the final variational question in
the Discussion of Section~7.2 of the source paper~\cite{CD}.

The answer is negative, independently of the choice of $E$.

\begin{theorem}[Universal factor-two obstruction]\label{thm:main}
Let $E\subset\R^2$ be measurable.  Suppose $\rho>0$ and $\omega$ are real
and sufficiently regular for the displayed expressions below, and suppose
$u=\rho e^{i\omega}\in\cS(\R^2)$.  Then
\begin{equation}\label{eq:factor2}
 \boxed{\quad
  2\int_E\Jac(\rho^2,\omega)
  \leq \int_{\R^2}|\nabla(\rho e^{i\omega})|^2.
 \quad}
\end{equation}
Consequently \eqref{eq:asked} is impossible for every $0<\delta\leq1$.
In particular, no set $E$ has the property requested in the source paper.
\end{theorem}

\section{Proof}

Put
\[
 J=\Jac(\rho^2,\omega).
\]
The polar decomposition gives the exact energy identity
\begin{equation}\label{eq:energy}
 |\nabla u|^2=|\nabla\rho|^2+\rho^2|\nabla\omega|^2.
\end{equation}
In particular $\rho\in L^2$, $\nabla\rho\in L^2$, and
$\rho\nabla\omega\in L^2$.  Since
\[
 J=2\rho(\rho_x\omega_y-\rho_y\omega_x),
\]
Cauchy--Schwarz and $2ab\leq a^2+b^2$ give the pointwise estimate
\begin{equation}\label{eq:pointwise}
 |J|
 \leq 2|\nabla\rho|\,\rho|\nabla\omega|
 \leq |\nabla\rho|^2+\rho^2|\nabla\omega|^2
 =|\nabla u|^2.
\end{equation}
Thus $J\in L^1(\R^2)$.

The global real phase is now decisive.  Define
\[
 V=(\rho^2\omega_y,-\rho^2\omega_x).
\]
Then, distributionally,
\begin{equation}\label{eq:div}
 \operatorname{div}V=J.
\end{equation}
Moreover
\[
 \|V\|_{L^1}
 \leq \|\rho\|_{L^2}\,\|\rho\nabla\omega\|_{L^2}<\infty.
\]
Let $\chi_R$ be a standard cutoff which is one on the ball $B_R$, vanishes
outside $B_{2R}$, and satisfies $|\nabla\chi_R|\leq C/R$.  From
\eqref{eq:div},
\[
 \left|\int_{\R^2}\chi_RJ\right|
 =\left|\int_{\R^2}\nabla\chi_R\cdot V\right|
 \leq \frac CR\|V\|_{L^1}\longrightarrow0.
\]
Dominated convergence, using $J\in L^1$, yields
\begin{equation}\label{eq:zero}
 \int_{\R^2}J=0.
\end{equation}

Write $J_+=\max\{J,0\}$ and $J_-=\max\{-J,0\}$.  Equation
\eqref{eq:zero} implies
\[
 \int J_+=\int J_-=\frac12\int|J|.
\]
For every measurable $E$ we therefore have, by \eqref{eq:pointwise},
\[
 \int_EJ
 \leq\int J_+
 =\frac12\int|J|
 \leq\frac12\int|\nabla u|^2,
\]
which is \eqref{eq:factor2}.

Finally, if \eqref{eq:asked} held for some $0<\delta\leq1$, its right-hand
side would be positive and \eqref{eq:factor2} would give
\[
 \int|\nabla u|^2
 <(1+\delta)\int_EJ
 \leq\frac{1+\delta}{2}\int|\nabla u|^2
 \leq\int|\nabla u|^2,
\]
a contradiction.  This proves the theorem. \qed

\section{Why the obstruction was hidden}

Pointwise near-equality in \eqref{eq:pointwise} is exactly the
Cauchy--Riemann mechanism suggested in the source.  Such a region creates
positive oriented Jacobian.  But a globally defined phase and Schwartz
decay force the total signed Jacobian to vanish by \eqref{eq:zero}.  Every
unit of positive Jacobian mass must therefore be balanced by an equal amount
of negative mass elsewhere, and \eqref{eq:pointwise} charges the energy of
both.  This doubles the lower bound and rules out constants approaching one.

The conclusion answers the displayed Schwartz-class problem exactly.  It
does not by itself characterize semigroup contractivity for every rough
matrix field, nor does it exclude variants in which the phase is not global
or the domain has a boundary carrying nonzero Jacobian flux.

\section{Novelty check and review note}

Cheap run indexes contained no result for arXiv:1611.00653.  Bounded searches
through 11 August 2026 using the exact variational wording, the paper title,
and the terms ``Schwartz'', ``Jacobian'', and ``global phase'' found the
source paper but no later answer to this question.  This is not an exhaustive
novelty claim.  The mathematical proof uses only the polar energy identity,
an $L^1$ divergence cutoff, and the positive/negative-part decomposition;
each integrability step is explicit above.

\begin{thebibliography}{9}
\bibitem{CD}
A. Carbonaro and O. Dragi\v cevi\'c,
\emph{Convexity of power functions and bilinear embedding for
divergence-form operators with complex coefficients},
J. Eur. Math. Soc. 22 (2020), 3175--3221,
\href{https://doi.org/10.4171/JEMS/984}{doi:10.4171/JEMS/984};
\href{https://arxiv.org/abs/1611.00653}{arXiv:1611.00653}.
\end{thebibliography}

\end{document}
